\section{Related Work}
\label{sec:related}

% Citation-safe Related Work.
% Verified BibTeX keys only: He2016DeepResidual, He2016IdentityMappings
% (arXiv API + Crossref DOIs; project README / paper_summary pointers).
% ResearchVault Analytical Notes: not available (Phase 2/3 waived_for_demo).
% Broader architecture--recipe literature beyond these two sources: [NEEDS_VERIFICATION].
% No results from this study; no invented citations.

\subsection{Residual Learning}
\label{sec:rw-residual}

Residual learning reformulates stacked layers as learning residual functions with respect
to layer inputs, rather than unreferenced mappings~\citep{He2016DeepResidual}.
Empirically, networks with identity shortcuts are reported to be easier to optimize than
plain counterparts of comparable depth, and residual nets are evaluated on large-scale
recognition benchmarks as well as on CIFAR-10 depth studies in the original
work~\citep{He2016DeepResidual}.
The CIFAR residual family uses a $6n{+}2$ stage layout with optional projection shortcuts
(Options~A/B); our study adopts this residual-learning framing as background, without
expanding into ImageNet-scale architectures beyond what that paper already reports.

\subsection{Pre-activation Residual Networks}
\label{sec:rw-preact}

A subsequent analysis of residual units emphasizes identity mappings on the skip path and
the role of after-addition activation in signal propagation~\citep{He2016IdentityMappings}.
That line of work motivates a \emph{pre-activation} residual unit in which batch
normalization and ReLU precede the convolutional transforms
($\mathrm{BN}{\rightarrow}\mathrm{ReLU}{\rightarrow}\mathrm{Conv}$ on the residual branch),
in contrast to the original \emph{post-activation} ordering where convolutions come first and
a nonlinearity is applied after the residual addition~\citep{He2016DeepResidual,He2016IdentityMappings}.
The pre-activation redesign is presented as improving trainability and generalization in the
settings studied by~\citet{He2016IdentityMappings}; we treat those claims as the source
hypothesis space for a controlled drop-in test, not as a literature consensus that every
post-activation CIFAR recipe must benefit from an untuned replacement.

\subsection{Architectural Variants and Training Recipes}
\label{sec:rw-recipes}

Both foundational papers introduce residual units together with concrete optimization
protocols (optimizer, schedule, data augmentation, and evaluation
settings)~\citep{He2016DeepResidual,He2016IdentityMappings}.
In particular,~\citet{He2016IdentityMappings} study how residual-block design choices
interact with forward/backward propagation and report ablation-style comparisons among unit
variants under their training setups.
This background supports reading block redesign and training procedure as coupled
experimental factors rather than as fully separable claims.
A broader survey of how later residual variants, normalizers, or modern training recipes
co-adapt with architecture is outside the verified bibliography of this draft
{[NEEDS\_VERIFICATION]}.

\subsection{Positioning of This Study}
\label{sec:rw-position}

This paper does not propose a new general-purpose network family.
It is a controlled empirical study on CIFAR-10 depth~20 with shortcut Option~A:
establish a frozen post-activation ResNet-20-A baseline under a fixed training recipe,
test a full pre-activation residual-unit replacement under that same recipe, and run
targeted ablations of the pre-activation BN+ReLU placement.
The research question and experimental boundary are therefore narrow by design---architecture
change under a frozen baseline protocol---and should not be read as a comprehensive
re-evaluation of pre-activation residual networks across datasets, depths, or retuned
schedules.
% End Related Work.
